\section{Top View of a Binary Tree}
\label{sec:trees-top-view}

\problemstatement{The top view is the set of nodes visible when the tree is
viewed from directly above: for each horizontal distance (HD), the single
\emph{topmost} node on that vertical line. Return these nodes ordered left
to right by HD.}

\begin{patternbox}
Reusing the HD tag from vertical order, \emph{``top view / viewed from
above''} means: on each vertical line, keep only the node encountered
\emph{first when scanning top-down}. A \textbf{BFS} scans strictly in
increasing depth, so the first node BFS places on a given HD is
automatically the topmost -- record an HD only the first time it appears.
\end{patternbox}

\subsection*{Understanding the problem carefully}

Looking down from above, a node is hidden if some other node shares its
horizontal distance but sits higher. Because BFS processes level $0$
before level $1$ before level $2$, the earliest time an HD is seen during
BFS corresponds to the highest node on that line. So we simply record the
first value per HD and ignore all later ones.

\begin{ideabox}[BFS, First Node per HD Wins]
Run level-order BFS carrying each node's HD (root $0$, left $-1$, right
$+1$). Keep a map from HD to value, writing an entry \emph{only if that HD
is not already present}. BFS's increasing-depth order guarantees the first
write for each HD is the topmost node. Emit the map in ascending HD order.
\end{ideabox}

\emph{Why BFS, not DFS:} a DFS might reach a deep node on some HD before a
shallower node on the same HD (down a different branch), wrongly recording
the deeper one. BFS's level-by-level order removes that hazard.

\subsection*{Algorithm}

\begin{enumerate}
  \item BFS from root with a queue of $(node, HD)$; root's HD $=0$.
  \item For each dequeued $(node, HD)$: if $HD$ not yet in the map, set
        $\textit{map}[HD]=node.val$. Enqueue $(left, HD-1)$ and
        $(right, HD+1)$.
  \item Output the map's values in ascending HD.
\end{enumerate}

\subsection*{Dry run}

\dryrun{Tree: root $1$; left $2$ (right child $4$); right $3$.}

\begin{itemize}
  \item $(1,0)$: map$\{0\!:\!1\}$; enqueue $(2,-1),(3,1)$.
  \item $(2,-1)$: map$\{-1\!:\!2,0\!:\!1\}$; enqueue $(4,0)$.
  \item $(3,1)$: map adds $1\!:\!3$.
  \item $(4,0)$: HD $0$ already present -- skip.
  \item Ascending HD: $2,1,3$.
\end{itemize}

\subsection*{C++ implementation}

\begin{lstlisting}
#include <bits/stdc++.h>
using namespace std;

struct TreeNode {
    int val;
    TreeNode* left;
    TreeNode* right;
    TreeNode(int v) : val(v), left(nullptr), right(nullptr) {}
};

vector<int> topView(TreeNode* root) {
    vector<int> ans;
    if (root == nullptr) return ans;

    map<int, int> topAtHD;               // hd -> first (topmost) value
    queue<pair<TreeNode*, int>> q;
    q.push({root, 0});
    while (!q.empty()) {
        auto [node, hd] = q.front();
        q.pop();
        if (topAtHD.find(hd) == topAtHD.end())
            topAtHD[hd] = node->val;     // first sighting wins
        if (node->left)  q.push({node->left,  hd - 1});
        if (node->right) q.push({node->right, hd + 1});
    }
    for (auto& [hd, v] : topAtHD) ans.push_back(v);
    return ans;
}

int main() {
    TreeNode* root = new TreeNode(1);
    root->left = new TreeNode(2);
    root->right = new TreeNode(3);
    root->left->right = new TreeNode(4);

    for (int x : topView(root)) cout << x << " ";
    cout << "\n";
    return 0;
}
\end{lstlisting}

\begin{complexitybox}
\textbf{Time Complexity.} $O(n\log n)$ -- $n$ BFS steps, each doing an
$O(\log n)$ map operation (or $O(n)$ overall with a hash map plus a final
sort). \textbf{Space Complexity.} $O(n)$ for the queue and map.
\end{complexitybox}

\begin{takeawaybox}
``First per HD under BFS'' is the whole trick for top view. The insistence
on BFS over DFS is the subtle correctness point: only a strictly
increasing-depth scan makes ``first seen'' equal to ``topmost.''
\end{takeawaybox}
